problem:  
Let \(M^6\) be a compact, simply connected Riemannian \(6\)-manifold with nonnegative sectional curvature admitting an effective, isometric \(T^3\)-action. The orbit space \(X=M/T^3\) is a compact \(3\)-dimensional Alexandrov space with nonnegative curvature whose stratification encodes orbit dimensions: open 3‑strata of principal orbits, 2‑strata corresponding to circle isotropy, and 1‑ or 0‑dimensional strata corresponding to \(T^2\) or \(T^3\) isotropy. Classical classification results recover \(M\) from the labeled orbit space \((X,\text{isotropy data})\), where isotropy weights determine how 2‑strata patch together.

**Weight‑Detection Problem:**  
For such \(M\), the metric link at a point \(p\in X\) is an Alexandrov space isometric to the orbit space of a linear \(T^3\)–representation at the corresponding orbit in \(M\). Each link carries metric data such as cone angles and curvature profiles that, in principle, depend on isotropy weights.  

Is the isotropy weight lattice (up to automorphism of \(T^3\)) *determined by the local Alexandrov geometry of links*?  

Equivalently:  
Let two \(6\)-manifolds \(M_1,M_2\) as above have orbit spaces \(X_1, X_2\) such that, around each singular point \(x_i\in X_i\), the space of directions is isometric to the same 2‑dimensional Alexandrov link. Must the corresponding isotropy weights at those orbits coincide up to lattice automorphism, and hence the labeled orbit spaces—and the homeomorphism type of \(M_i\)—agree?  

In other words, can the quantitative metric data of the Alexanderov quotient (cone angles, link curvatures) reveal the full isotropy representation, thereby making the labeling information metric in nature?

Why is it a "good" problem:  
This formulation directly targets the core geometric mystery underlying the metric classification of nonnegatively curved torus manifolds: *Is isotropy algebraic data recoverable from curvature‑constrained local metrics?* It is a sharpened, local version of the reconstruction problem, focusing on the precise analytic-geometric interface between representation theory and Alexandrov geometry.  

Resolving it would illuminate whether the “nonnegatively curved orbit geometry” already encodes group representation data geometrically via its links and angles, or whether isotropy information is inherently non‑metric, requiring discrete combinatorial labels. Either outcome reshapes understanding of how curvature interacts with symmetry.  

The question is extremely simple to state but conceptually deep—it asks whether “metric spaces of directions know their isotropy weights.” It bridges Riemannian curvature bounds, Alexandrov link geometry, and torus representation theory, thus meeting the criteria for a profoundly integrative and foresighted research problem.